% gerado por scripts/revisao/build_revisao_tabs.py -- nao editar a mao
\begin{table}[htbp]
\centering
\small
\setlength{\tabcolsep}{5pt}
\caption{Whether the temperature gain of Table~\ref{tab:temperatura} survives once the
boosting models are given the annual cycle explicitly. Gain is positive when the covariate
helps.}
\label{tab:tempdiffcal}
\begin{tabular}{lrrr}
\toprule
Model and feature set & Without temp. & With temp. & Gain (pp) \\
\midrule
CatBoost, lags only & 6.59 & 6.33 & $+0.258$ \\
CatBoost, + calendar and seasonal difference & 6.13 & 6.05 & $+0.076$ \\
\midrule
XGBoost, lags only & 6.83 & 6.51 & $+0.327$ \\
XGBoost, + calendar and seasonal difference & 6.35 & 6.39 & $-0.036$ \\
\bottomrule
\end{tabular}

\vspace{0.5em}
\begin{minipage}{\textwidth}
\footnotesize
The first row of each block is the specification of Table~\ref{tab:temperatura}, lagged
values only. The second adds Fourier calendar terms and takes the seasonal difference
$y_t - y_{t-12}$ as the target, so the model no longer has to infer the annual cycle from
the covariate. The gain shrinks in both models, CatBoost from $+0.258$ to $+0.076$ and XGBoost from $+0.327$ to $-0.036$ percentage
points, and changes sign for XGBoost. The reading is that a large part of what minimum
temperature contributed in Table~\ref{tab:temperatura} was the month of the year rather
than the weather, and that the remaining contribution is too small to separate from noise.
The comparison that carries this conclusion is within a row, the same run with and without
the covariate; the without-temperature column of the XGBoost blocks is not identical to
Table~\ref{tab:temperatura} because XGBoost does not reproduce across library versions,
a limitation already declared in the manuscript.
\end{minipage}
\end{table}